% ============================================================================
% Fichier  : partie3/P03_C08_loi_cadre_africain.tex
% Titre    : La Loi n° 2024/017 et le Cadre Africain
% Auteur   : MINKA MI NGUIDJOÏ Thierry Emmanuel
% Version  : 1.0 -- novembre 2026
% Corrigés : corriges/Corr_C08.tex
% ============================================================================
% Sources : loi n° 2024/017 (texte officiel, titres I à VII) ; loi
% n° 2010/012, art. 41 à 47 (texte officiel) ; Convention de Malabo ;
% Convention de Budapest ; acte additionnel CEDEAO A/SA.1/01/10 ; RGPD.
% Statut des conventions d'après Minka Mi Nguidjoï, Massoussi Songo et
% Bialo (2026, manuscrit) : Budapest, Cameroun Partie le 15/12/2023, en
% vigueur le 01/04/2024 ; Malabo, signée le 12/08/2021, non ratifiée.
% ============================================================================

\chapter{La Loi n° 2024/017 et le Cadre Africain}
\label{chap:loi}
\chaptermark{La loi et le cadre africain}

\epigraph{%
    « On cite une loi pour ce qu'elle dit,\\
    pas pour ce que l'on aimerait qu'elle dise. »%
}{Principe du cours}

\niveauChapitre{Fondamental}{%
    Chapitres~\ref{chap:donnees} à~\ref{chap:securite}~; droit du numérique
    (UE MSN-2S3-JUR01)
}

\begin{boxobjectifs}
\begin{itemize}
    \item Présenter l'architecture de la loi n°~2024/017 en sept titres.
    \item Expliquer le régime d'autorisation préalable et ses conséquences
          tant que l'Autorité n'est pas organisée.
    \item Décrire les missions de l'Autorité et l'échelle des sanctions
          administratives, civiles et pénales.
    \item Articuler la loi n°~2024/017 avec la loi n°~2010/012.
    \item Situer la loi par rapport à la Convention de Malabo, à la
          Convention de Budapest, au droit de la CEDEAO et au RGPD.
\end{itemize}
\end{boxobjectifs}

\bigskip

% ============================================================================
\section*{Avant de commencer : « c'est le RGPD camerounais »}
% ============================================================================

Bafoussam, lundi 3~août~2026. Le directeur général de NYANGA CRÉDIT a reçu
d'un cabinet de conseil une offre de « mise en conformité RGPD
camerounais » en douze semaines. Le document promet la désignation d'un
délégué à la protection des données « obligatoire », une procédure de
notification des violations « sous 72 heures », et la constitution d'un
dossier d'« intérêt légitime » pour la vidéosurveillance.

M.~NGANDO, juriste de KORO AUDIT, pose la proposition à côté du texte de la
loi. Aucune de ces trois affirmations ne s'y trouve. La loi ne connaît pas
de délégué obligatoire, elle dit « sans délai » et non 72~heures, et elle
ignore l'intérêt légitime. À l'inverse, l'offre ne dit pas un mot de
l'autorisation préalable de l'article~19, qui est l'obligation la plus
lourde de la loi.

\begin{boxalert}
\textbf{Violation typique.} Lire la loi camerounaise à travers le RGPD, et
lui attribuer des obligations européennes qu'elle ne contient pas, en
oubliant les siennes.

\smallskip
\textbf{Règle~:} l'audit applique le texte officiel. Le droit européen sert
de comparaison et de source de bonnes pratiques, déclarées comme telles~;
il ne crée ni obligation ni dispense au Cameroun, sauf pour les traitements
qu'il régit lui-même.
\end{boxalert}

% ============================================================================
\section{L'Architecture de la Loi}
\label{sec:c08-architecture}
% ============================================================================

\begin{tableaucours}{cL{5.6cm}cL{4.0cm}}{%
    \tableauHeader{Titre} &
    \tableauHeader{Objet} &
    \tableauHeader{Art.} &
    \tableauHeader{Chapitres du cours}
}
I & Dispositions générales~: objet, champ, définitions & 1--5 &
\ref{chap:donnees} \\
II & Régime juridique~: principes, formalités préalables, obligations,
interconnexion & 6--36 & \ref{chap:principes}, \ref{chap:cartographie},
\ref{chap:risque}, \ref{chap:securite}, \ref{chap:soustraitants},
\ref{chap:transferts}, \ref{chap:violations} \\
III & Droits de la personne concernée & 37--47 & \ref{chap:droits},
\ref{chap:profilage} \\
IV & Interdictions & 48--52 & \ref{chap:donnees}, \ref{chap:secteurs} \\
V & Autorité de protection & 53 & ce chapitre \\
VI & Sanctions administratives, civiles et pénales & 54--71 & ce chapitre \\
VII & Dispositions diverses, transitoires et finales & 72--75 &
\ref{chap:prologue}, \ref{chap:secteurs} \\
\end{tableaucours}
\vspace{-0.8em}
\begin{center}{\small\itshape Tableau~8.1. Les sept titres de la loi n°~2024/017}\end{center}

La loi compte soixante-quinze articles. Elle « abroge toutes les
dispositions antérieures contraires » (article~75) et renvoie à des textes
particuliers les modalités d'application (article~74). Une quinzaine de ses
dispositions dépendent d'un texte réglementaire ou d'un référentiel de
l'Autorité~: critères de risque, évaluation d'impact, durées de
conservation, délais des droits, modalités d'autorisation et de contrôle.

% ============================================================================
\section{Le Régime d'Autorisation Préalable}
\label{sec:c08-autorisation}
% ============================================================================

\begin{boxjuris}
\textbf{\loicam{n°~2024/017}, article~19}~\cite{loi2024017}~:
\textit{« (1) Le traitement des données à caractère personnel est soumis à
l'obtention préalable d'une autorisation délivrée par l'Autorité de
protection des données à caractère personnel. (2) Tout processus
d'interconnexion et d'interopérabilité des fichiers de données sensibles,
relatives aux mineurs, est soumis à l'obtention préalable d'une
autorisation délivrée par l'Autorité de protection des données à caractère
personnel. (3) Les modalités de délivrance de l'autorisation visée aux
alinéas 1 et 2 ci-dessus sont fixées par voie réglementaire. »}

\smallskip
\textbf{Portée.} La loi ne prévoit ni déclaration simplifiée ni dispense~:
tout traitement est soumis à autorisation. L'article~49 interdit le
traitement sans autorisation préalable, et l'article~55 le punit d'une
pénalité de 5 à 50~millions de francs CFA.
\end{boxjuris}

Le régime est l'inverse du modèle européen, qui a supprimé les formalités
préalables au profit de la responsabilisation du responsable. Il rapproche
la loi camerounaise des premières lois africaines et françaises, fondées
sur un contrôle \textit{a priori}.

Ce régime pose une difficulté pratique~: l'autorisation est délivrée par une
Autorité dont « la création, ainsi que l'organisation et le
fonctionnement » sont fixés par décret (article~53, alinéa~2), et selon des
modalités elles-mêmes réglementaires (article~19, alinéa~3). Tant que ces
textes ne sont pas publiés et que l'Autorité ne reçoit pas de demandes,
l'organisation ne peut pas obtenir d'autorisation. Que doit coter
l'auditeur~?

\begin{itemize}
    \item « Le traitement est autorisé »~: \NON{}, puisqu'aucune autorisation
          n'existe. La cote est exacte~; elle ne dit pas que l'écart est
          imputable à l'organisation.
    \item « L'organisation a fait ce qui dépendait d'elle »~: \OUI{} si elle
          a préparé le dossier (registre, évaluations, mesures) et peut
          prouver qu'elle le déposera dès que la procédure sera ouverte~;
          \NON{} sinon.
\end{itemize}

Le rapport présente les deux cotes ensemble. À NYANGA, la seconde est
\NON{}~: aucun dossier n'a été préparé.

% ============================================================================
\section{L'Autorité de Protection}
\label{sec:c08-autorite}
% ============================================================================

\begin{boxjuris}
\textbf{\loicam{n°~2024/017}, article~53, alinéa~1 (extraits)}~\cite{loi2024017}~:
\textit{« L'Autorité de protection des données à caractère personnel est un
organisme public indépendant chargé, notamment~: de veiller à l'application
de la présente loi, des textes pris pour son application ainsi que des
conventions internationales en la matière~; de délivrer les autorisations
prévues par la présente loi, assorties des cahiers de charges
conséquents~; d'élaborer, publier et mettre à jour le référentiel des
mesures techniques et organisationnelles des données à caractère
personnel~; [\ldots] de traiter les réclamations, pétitions et plaintes
[\ldots]~; d'établir et publier la liste des pays reconnus comme offrant un
niveau de protection des données à caractère personnel équivalent aux
exigences prévues par le droit camerounais~; [\ldots] »}
\end{boxjuris}

Trois instruments de l'Autorité conditionnent l'audit~: le
\textbf{référentiel des mesures} (articles~24, 27, 28, 57), les
\textbf{cahiers des charges} attachés aux autorisations (article~61), et la
\textbf{liste des pays} offrant une protection équivalente
(Chapitre~\ref{chap:transferts}). Tant qu'ils n'existent pas, les exigences
qui en dépendent sont des \NEK{} de critère.

\begin{boxinfo}
\textbf{État de l'Autorité au moment de la rédaction.} À la date de ce cours
(novembre~2026), l'auteur n'a pas trouvé de texte publié organisant
l'Autorité, et ses coauteurs et lui relèvent dans leurs travaux en cours que
ce texte était attendu après l'échéance du 23~juin~2026. Le lecteur
vérifiera la situation au moment où il lit~: si l'Autorité est organisée,
une partie des \NEK{} de critère de ce livre se lève.
\end{boxinfo}

% ============================================================================
\section{Les Sanctions}
\label{sec:c08-sanctions}
% ============================================================================

\subsection{Sanctions administratives}

\begin{boxjuris}
\textbf{\loicam{n°~2024/017}, article~54}~\cite{loi2024017}~:
\textit{« (1) Lorsque l'Autorité de protection des données à caractère
personnel constate le non-respect des obligations mises à la charge du
responsable du traitement ou de son sous-traitant, celui-ci est mis en
demeure de s'y conformer dans un délai maximum de dix (10) jours. (2) Passé
le délai prévu à l'alinéa 1 ci-dessus, l'Autorité de protection des données
à caractère personnel adresse une injonction de mise en conformité du
traitement, sous astreinte dont le montant ne peut excéder cent mille
(100~000) francs CFA par jour de retard. [\ldots] »}

\smallskip
\textbf{Portée.} À défaut de mise en conformité, l'Autorité peut suspendre
l'activité, retirer l'autorisation ou interdire tout traitement
(alinéa~3).
\end{boxjuris}

\begin{tableaucours}{cL{6.4cm}L{4.0cm}}{%
    \tableauHeader{Art.} &
    \tableauHeader{Manquement} &
    \tableauHeader{Montant (F CFA)}
}
55 & Traitement sans autorisation préalable & 5 à 50~millions \\
56 & Refus de communiquer à la personne les informations la concernant & 1
à 10~millions \\
57 & Non-respect du référentiel & 5 à 20~millions \\
58 & Non-respect des règles d'interconnexion & 1 à 5~millions \\
59 & Traitement sans la certification prévue & 5 à 15~millions \\
60 & Transfert vers un pays tiers sans autorisation & 10 à 50~millions \\
61 & Non-respect du cahier des charges & 10 à 100~millions \\
\end{tableaucours}
\vspace{-0.8em}
\begin{center}{\small\itshape Tableau~8.2. Pénalités et amendes administratives (articles~55 à~61)}\end{center}

\subsection{Sanctions civiles et pénales}

L'article~62 ouvre à la personne concernée le référé « en cas d'atteinte
grave aux droits », éventuellement sous astreinte, et l'action en
réparation. Les articles~63 à~70 créent des délits~; l'article~71 rend les
personnes morales pénalement responsables, avec une amende de 50~millions à
un milliard de francs CFA.

\begin{tableaucours}{cL{6.6cm}L{3.8cm}}{%
    \tableauHeader{Art.} &
    \tableauHeader{Délit} &
    \tableauHeader{Peine d'emprisonnement}
}
63 & Collecte ou accès frauduleux, déloyal ou illicite (peines doublées
avec verrouillage ou cryptage) & 2 à 5~ans \\
64 & Traitement maintenu malgré l'opposition, refus de rectifier ou
d'effacer & 1 à 3~ans \\
65 & Traitement à des fins de profilage & 3 à 10~ans \\
66 & Données sensibles sans consentement exprès~; récupération
frauduleuse de données supprimées & 6~mois à 2~ans \\
67 & Détournement de finalité, traitement ultérieur incompatible & 6~mois
à 2~ans \\
68 & Divulgation portant atteinte à la dignité ou à la vie privée (sur
plainte) & 6~mois à 2~ans \\
69 & Transfert vers l'étranger en méconnaissance de la loi & 3 à 10~ans \\
70 & Entrave à l'action de l'Autorité & 6~mois à 3~ans \\
\end{tableaucours}
\vspace{-0.8em}
\begin{center}{\small\itshape Tableau~8.3. Les délits de la loi (amendes assorties, non reproduites)}\end{center}

L'article~65 mérite attention~: il punit « le responsable de traitement ou
le sous-traitant qui procède ou fait procéder à un traitement de données à
caractère personnel à des fins de profilage », sans condition apparente.
Sa portée exacte, pour une notation de crédit par exemple, est discutée au
Chapitre~\ref{chap:profilage}.

% ============================================================================
\section{La Loi n° 2010/012, Toujours Applicable}
\label{sec:c08-loi2010}
% ============================================================================

La loi n°~2010/012 relative à la cybersécurité et à la cybercriminalité
contient une section consacrée à la vie privée (articles~41 à~47), qui
reste applicable en ce qu'elle n'est pas contraire à la loi de 2024.

\begin{boxjuris}
\textbf{\loicam{n°~2010/012}, article~47}~\cite{loi2010012}~:
\textit{« L'utilisation des réseaux de communications électroniques et des
systèmes d'information aux fins de stocker les informations ou d'accéder à
des informations stockées dans un équipement terminal d'une personne
physique ou morale, ne peut se faire qu'avec son consentement
préalable. »}

\smallskip
\textbf{Portée.} Un kit de mesure d'audience intégré à une application
mobile, qui lit l'identifiant du téléphone, accède à des informations
stockées dans l'équipement terminal~: il exige un consentement préalable.
C'est le cas du traitement découvert dans l'application de NYANGA
(Chapitre~\ref{chap:cartographie}).
\end{boxjuris}

L'article~44 de la même loi interdit d'écouter, d'intercepter ou de stocker
les communications « sans le consentement des utilisateurs concernés, sauf
lorsque cette personne y est légalement autorisée », et l'article~45
autorise l'enregistrement des communications « effectué dans le cadre
professionnel en vue de fournir la preuve numérique d'une communication
électronique ». L'enregistrement des appels du service client de NYANGA
peut s'appuyer sur l'article~45 s'il sert à prouver les opérations
convenues par téléphone~; il ne le peut pas s'il sert à évaluer les agents
ou à entraîner un logiciel. La finalité décide.

% ============================================================================
\section{Le Cadre Africain et International}
\label{sec:c08-international}
% ============================================================================

\begin{tableaucours}{L{3.4cm}L{4.2cm}L{4.6cm}}{%
    \tableauHeader{Instrument} &
    \tableauHeader{Statut pour le Cameroun} &
    \tableauHeader{Usage dans l'audit}
}
Convention de Malabo (UA, 2014) & Signée le 12~août~2021, non ratifiée
(d'après~\cite{minka2026chainepenale}) & Référence comparative et
prospective, sans obligation opposable \\
Convention de Budapest (2001) & Partie depuis le 15~décembre~2023, en
vigueur le 1\textsuperscript{er}~avril~2024 (d'après~\cite{minka2026chainepenale})
& Coopération pénale~; tension avec les transferts
(Chapitre~\ref{chap:recherche}) \\
Acte additionnel CEDEAO (2010) & Non applicable (Cameroun hors CEDEAO) &
Comparaison régionale \\
RGPD (UE, 2016) & Non applicable en tant que tel & Comparaison~; applicable
aux prestataires établis dans l'Union \\
\end{tableaucours}
\vspace{-0.8em}
\begin{center}{\small\itshape Tableau~8.4. Les instruments internationaux et leur statut}\end{center}

La Convention de Malabo~\cite{ua2014malabo}, entrée en vigueur le 8~juin~2023,
consacre son chapitre~II à la protection des données à caractère
personnel. Faute de ratification, elle n'oblige pas le Cameroun~; elle
éclaire cependant les choix de la loi, et elle pourrait devenir applicable.
L'auditeur la cite comme référentiel comparatif, jamais comme critère
d'audit.

Le RGPD~\cite{ue2016rgpd} ne s'applique pas aux traitements de NYANGA
CRÉDIT. Il s'applique en revanche à un prestataire de nuage établi dans
l'Union européenne pour ses propres obligations de sous-traitant~: NOUN PAY
le rencontrera dans les contrats de son hébergeur
(Chapitre~\ref{chap:transferts}).

\begin{tableaucours}{L{3.6cm}L{4.4cm}L{4.2cm}}{%
    \tableauHeader{Point} &
    \tableauHeader{Loi n°~2024/017} &
    \tableauHeader{RGPD}
}
Formalité préalable & Autorisation pour tout traitement & Aucune (sauf
consultation préalable) \\
Bases du traitement & Consentement et trois dérogations & Six bases, dont
contrat et intérêt légitime \\
Délégué à la protection & Non prévu & Obligatoire dans certains cas \\
Notification d'une violation & « Sans délai », à l'Autorité et à la
personne & 72~heures à l'autorité~; personnes si risque élevé \\
Données sensibles & Liste ouverte, dont transactions bancaires et origine
régionale & Liste fermée \\
Profilage & Délit (article~65) & Encadré, non interdit \\
Sanctions maximales & 1~milliard F CFA (personne morale, pénal) & 4~\% du
chiffre d'affaires mondial \\
\end{tableaucours}
\vspace{-0.8em}
\begin{center}{\small\itshape Tableau~8.5. Sept différences avec le RGPD}\end{center}

% ============================================================================
\section{L'Exposition de NYANGA}
\label{sec:c08-nyanga}
% ============================================================================

\begin{boxscenario}{Mission NYANGA. La note juridique du 10 août 2026}
\textbf{Constatations} (rapprochement des constats d'audit et des
articles). Traitements sans autorisation~: 47 (article~55). Demandes
d'accès sans réponse~: 12 (article~56). Oppositions non exécutées~: 3
(article~64). Interconnexion sans obligation légale~: 1, le classeur des
mauvais payeurs (article~58). Données sensibles sans consentement exprès~:
au moins 4 traitements (article~66). Kit de mesure d'audience sans
consentement~: 1 (loi n°~2010/012, article~47).

\textbf{Analyse} (confiance moyenne). L'exposition théorique est élevée,
mais sa réalisation dépend de l'organisation de l'Autorité, de l'action
des personnes et du parquet. M.~NGANDO déconseille de chiffrer une
« amende probable »~: aucune décision n'existe encore pour en fonder une
estimation. Il classe les écarts par gravité pour les personnes, pas par
montant.
\end{boxscenario}

% ============================================================================
\section{Ce que la Loi demande à chaque Métier}
\label{sec:c08-metiers}
% ============================================================================

\begin{boxmetier}{Juriste conformité}
Lire le texte officiel, distinguer les obligations applicables, les
obligations suspendues à un texte d'application et les bonnes pratiques
importées, et le dire dans chaque avis.
\end{boxmetier}

\begin{boxmetier}{Auditeur de la protection des données}
Coter séparément « l'exigence est satisfaite » et « l'organisation a fait ce
qui dépendait d'elle » quand l'écart tient à l'absence de l'Autorité.
\end{boxmetier}

\begin{boxmetier}{Responsable de traitement}
Préparer dès maintenant les dossiers d'autorisation~: registre, fondements,
évaluations, mesures. Le jour où la procédure ouvrira, ce sont eux qui
seront demandés.
\end{boxmetier}

\begin{boxresume}
\begin{itemize}
    \item La loi compte sept titres et soixante-quinze articles~; une
          quinzaine de dispositions attendent un texte d'application.
    \item Tout traitement est soumis à autorisation préalable de
          l'Autorité (article~19)~; tant qu'elle n'est pas organisée,
          l'auditeur cote aussi ce qui dépendait de l'organisation.
    \item Les sanctions vont de la mise en demeure à dix jours jusqu'à un
          milliard de francs CFA pour une personne morale.
    \item La loi n°~2010/012 reste applicable, notamment son article~47 sur
          l'accès aux équipements terminaux.
    \item La Convention de Malabo, signée mais non ratifiée, est une
          référence comparative~; le RGPD n'oblige pas les organisations
          camerounaises en tant que telles.
\end{itemize}
\end{boxresume}

\begin{boxinfo}
\textbf{Points de contrôle.} Vous avez compris ce chapitre si vous savez~:
\begin{itemize}
    \item citer trois différences entre la loi et le RGPD~;
    \item expliquer comment coter l'absence d'autorisation quand l'Autorité
          n'est pas organisée~;
    \item rattacher un écart à l'article qui le sanctionne.
\end{itemize}
\end{boxinfo}

\begin{boxexo}[title={\ding{45}\quad Exercice~8.1. Quel article \niveaubase}]
Indiquez l'article qui sanctionne~: (a) un transfert vers un hébergeur
étranger sans autorisation~; (b) un agent qui récupère des fichiers
effacés de la corbeille d'un collègue~; (c) un établissement qui refuse de
communiquer à un client ses données~; (d) une entreprise qui refuse de
recevoir les agents de l'Autorité.
\end{boxexo}

\begin{boxexo}[title={\ding{45}\quad Exercice~8.2. L'offre du cabinet \niveaumoyen}]
Reprenez l'offre de « mise en conformité RGPD camerounais » reçue par
NYANGA. Corrigez ses trois affirmations et ajoutez les trois obligations
camerounaises qu'elle omet, articles à l'appui.
\end{boxexo}

\begin{boxexo}[title={\ding{45}\quad Exercice~8.3. Les appels enregistrés \niveaumoyen}]
Le service client enregistre tous les appels, les conserve deux ans et les
écoute pour évaluer les agents. Analysez ce traitement au regard des
articles~44 et~45 de la loi n°~2010/012 et des articles~9 et~11 de la loi
n°~2024/017. Proposez une règle conforme.
\end{boxexo}

\begin{boxexo}[title={\ding{45}\quad Exercice~8.4. Sans Autorité \niveauexpert}]
Un client de MBOLO dépose plainte auprès de l'Autorité en septembre~2026.
L'Autorité n'est pas encore organisée. (a) Quelles autres voies s'offrent
à lui~? (b) Quel est l'intérêt, pour MBOLO, d'avoir préparé ses dossiers
d'autorisation~? (c) Comment l'auditeur rend-il compte de cette
situation dans son rapport~?
\end{boxexo}

\bigskip

\begin{boxinfo}
\textbf{Pour aller plus loin.}
\begin{itemize}
    \item La loi n°~2024/017~\cite{loi2024017} et la loi
          n°~2010/012~\cite{loi2010012}, dans leur texte officiel.
    \item La Convention de Malabo~\cite{ua2014malabo}, chapitre~II.
    \item L'acte additionnel de la CEDEAO~\cite{cedeao2010acte}, pour une
          comparaison régionale.
\end{itemize}
\end{boxinfo}

\begin{boxinfo}
\textbf{Transition.} La loi fixe les obligations~; les normes fournissent la
manière de les organiser et de les certifier. L'auditeur, lui, obéit à une
déontologie propre. C'est l'objet du Chapitre~\ref{chap:normes}.
\end{boxinfo}
